\chapter{Конструкторская часть}

В конструкторской части будет разработан комбинированный метод решения задачи коммивояжёра, основанный на муравьином алгоритме с динамическим подбором гетерогенных параметров с использованием генетического алгоритма. Будут описаны общая структура метода, детальные этапы его работы, используемые структуры данных, а также архитектурное разбиение на модули. Для наглядности основные алгоритмические решения будут представлены в виде IDEF0-диаграмм и схем алгоритмов.

\section{Обобщённый генетический алгоритм}

Для разработки предлагаемого комбинированного метода, необходимо рассмотреть основные принципы генетического алгоритма. Это обусловлено тем, что предлагаемый метод использует генетический алгоритм как механизм адаптации гетерогенных параметров муравьёв.

Генетический алгоритм представляет собой метаэвристику, вдохновлённую принципами естественного отбора и эволюции в живой природе. Впервые подобный подход был предложен Джоном Холландом в 1975 году~\cite{genetic}. Основная идея заключается в моделировании биологического процесса эволюции: популяция потенциальных решений задачи (особей) развивается из поколения в поколение за счёт применения операторов отбора, скрещивания и мутации. При этом более приспособленные особи, то есть дающие лучшее значение целевой функции, имеют больше шансов передать свои характеристики потомкам. В результате последовательных итераций популяция в целом улучшается, приближаясь к оптимальному решению.

\subsection{Понятия}

В рамках генетического алгоритма используются следующие базовые понятия~\cite{genetic-engineer}:
\begin{itemize}
	\item Особь -- потенциальное решение задачи, представленное. Каждая особь обладает определённым значением приспособленности;
	\item Хромосома -- структура данных, кодирующая решение. Она состоит из генов, расположенных в определённом порядке. В зависимости от задачи гены могут быть бинарными, целочисленными, вещественными или иметь другую природу;
	\item Ген -- элементарная единица хромосомы, отвечающая за кодирование отдельного параметра решения. Совокупность генов полностью определяет характеристики особи;
	\item Аллель -- конкретное значение гена;
	\item Популяция -- множество особей, существующих в рамках одного поколения. Размер популяции является фиксированным параметром и определяет количество одновременно рассматриваемых решений;	
	\item Поколение -- одна итерация эволюционного процесса, в результате которой из текущей популяции формируется новая;	
	\item Фитнес-функция (функция приспособленности) -- критерий качества особи, который оценивает, насколько хорошо данное решение удовлетворяет поставленной задаче. Значение фитнес-функции определяет вероятность участия особи в размножении;
	\item Операторы селекции, скрещивания, мутации -- операции применяемые к хромосомам для получения и/или изменения особей. 
\end{itemize}


\subsection{Схема генетического алгоритма}

Схема классического генетического алгоритма представлена на рисунке~\ref{fig:genetic}.

\begin{figure}[H]
	\centering
	\includegraphics[width=0.6\textwidth]{genetic/genetic}
	\caption{Схема генетического алгоритма}
	\label{fig:genetic}
\end{figure}


\textbf{Инициализация популяции}

Формируется начальная популяция заданного размера. Каждая особь создаётся случайным образом или некоторым заданным алгоритмом.

\textbf{Оценка приспособленности}

Для каждой особи вычисляется значение фитнес-функции. Полученные значения позволяют ранжировать особей по качеству и используются на этапе селекции.

\textbf{Селекция (отбор)}

Из текущей популяции выбираются особи, которые станут родителями для создания потомства.
Существуют различные алгоритмы селекции, отличающиеся случайностью и стойкостью выбора родителей. Алгоритм обычно устроен таким образом, чтобы лучшие особи имели больше шансов стать родителями, но при этом сохранялось определённое разнообразие для предотвращения преждевременной сходимости.

\textbf{Кроссинговер (скрещивание)}

Из отобранных родителей формируются пары, которые в результате применения оператора кроссинговера создают одну или две особи-потомка. В зависимости от реализации оператора может происходить обмен участками хромосом (аллелями), их математическая комбинация или иное преобразование, наследующее признаки родителей.

\textbf{Мутация}

Каждый ген полученных потомков с малой вероятностью подвергается случайному изменению. Мутация поддерживает генетическое разнообразие популяции и позволяет исследовать новые области пространства поиска, которые могли быть утеряны в процессе селекции и кроссинговера

\textbf{Проверка критерия остановки}

Эволюционный процесс повторяется до выполнения одного из критериев:
\begin{itemize}
	\item достижение максимального числа поколений;
	\item достижение требуемого значения фитнес-функции;
	\item отсутствие улучшений в течение заданного числа поколений (стагнация);
	\item превышение допустимого времени работы.
\end{itemize}

\textbf{Таким образом}, генетический алгоритм оперирует популяцией особей, каждая из которых кодирует потенциальное решение. Последовательное применение операторов селекции, кроссинговера и мутации позволяет эволюционировать популяции в направлении улучшения качества решений. В предлагаемом комбинированном методе генетический алгоритм планируется к использованию для адаптации параметров $\alpha$ и $\beta$ муравьёв.

\section{Описание предлагаемого метода}

\subsection{Требования к разрабатываемому методу}

К разрабатываемому методу предъявляются следующие функциональные требования:
\begin{itemize}
	\item возможность решения задачи коммивояжёра для полных взвешенных неориентированных графов;
	\item адаптация параметров муравьёв в процессе поиска.
\end{itemize}

и нефункциональные требования:
\begin{itemize}
	\item улучшение сходимости по сравнению с классическим муравьиным алгоритмом;
	\item лучшая устойчивость к выбору начальных параметров $\alpha$ и $\beta$ по сравнению с классическим муравьиным алгоритмом.
\end{itemize}

\subsection{Формализация метода}

На рисунке~\ref{fig:02A0} представлено верхнеуровневое разбиение метода на последовательность действий. Блоки A1-A3 являются частями муравьиного алгоритма, а блок A4 описывает динамическую адаптацию параметров муравьев.

\begin{figure}[H]
	\centering
	\includegraphics[width=\textwidth]{method/02\_A0}
	\caption{IDEF0 диаграмма метода 1-го уровня}
	\label{fig:02A0}
\end{figure}

\begin{figure}[H]
	\centering
	\includegraphics[width=\textwidth]{method/03\_A2}
	\caption{IDEF0 диаграмма формирования путей муравьёв}
	\label{fig:03A2}
\end{figure}

Диаграмма декомпозиции блока A2 (рисунок~\ref{fig:03A2}) описывает процесс формирования путей муравьями, совпадающей с классическим муравьиным алгоритмом, с добавлением локальной оптимизации для избранных путей поколения.

\begin{figure}[H]
	\centering
	\includegraphics[width=\textwidth]{method/03\_A4}
	\caption{IDEF0 диаграмма формирования новой популяции}
	\label{fig:03A4}
\end{figure}

Диаграмма декомпозиции блока A4 (рисунок~\ref{fig:03A4}) описывает работу генетического алгоритма со всеми его классическими этапами, однако особями в данном методе являются не решения задачи, а муравьи, в частности два их индивидуальных параметра: $\alpha$ и $\beta$.

Таким образом, генетический алгоритм интегрируется как метаэвристика над муравьиным алгоритмом, которая динамически адаптирует параметры муравьев под конфигурацию графа и текущего уровня феромонов.

\section{Поиск муравьиным алгоритмом}

На основе представленной IDEF0-формализации и выделенных параметров метода ниже приводится детальное описание каждого этапа. Для каждого этапа описываются используемые алгоритмы, формулы и параметры, а также их обоснование с точки зрения цели разрабатываемого комбинированного подхода.

\subsection{Инициализация первого поколения}

\textbf{Инициализация феромонов}

Так же как и в классическом муравьином алгоритме на каждом ребре устанавливается в небольшое положительное значение уровня феромона $\tau_{ij}(0)=\tau_0$.

Это позволяет сгладить первую итерацию, увеличивая шансы на выбор муравьём более длинных участков графа, что увеличивает начальную область поиска и разнообразие решений.

\textbf{Формирование начальной популяции муравьёв}

Создаётся популяция размером $m$, где $m$ -- количество муравьёв в одной популяции. Каждый муравей имеет индивидуальные, но пока что идентичные $\alpha=\alpha_0$ и $\beta=\beta_0$.

\subsection{Формирование путей муравьиным алгоритмом}

Данный этап выполняется на каждой итерации для всех муравьёв, в результате чего находится множество из $m$ решений.

\textbf{Выбор начальной вершины} — для каждого муравья выбирается случайным образом, что снижает влияние начального выбора на результат.

\textbf{Правило выбора очередной вершины} -- вероятностное, аналогично классическому с поправкой на гетерогенность параметров: Пусть на итерации $k$ муравей $a$ находится в вершине $v_i$, а множество ещё не посещённых им вершин обозначается как $U_a(i)$, тогда вероятность выбора вершины $v_j \in U_a(i)$ в качестве следующей задаётся формулой~(\ref{eq:cprobrule}).

\begin{equation}
	P_{ij}^{(a)}(k) =
	\begin{cases}
		\displaystyle \frac{\tau_{ij}(k)^{\alpha_a} \cdot \eta_{ij}^{\beta_a}}
		{\sum\limits_{v_l \in U_a(i)} \tau_{il}(k)^{\alpha_a} \cdot \eta_{il}^{\beta_a}}, & \text{если } v_j \in U_a(i); \\
		0, & \text{иначе}, \\
	\end{cases}
	\label{eq:cprobrule}
\end{equation}
где
\begin{itemize}
	\item $\tau_{ij}(k)$ -- интенсивность феромона на ребре $e_{ij}$ на итерации $k$;
	\item $\eta_{ij} = \frac{1}{F(e_{ij})}$ -- эвристическая значимость ребра (обратная длина);
	\item $\alpha_a$ и $\beta_a$ -- индивидуальные параметры муравья $a$.
\end{itemize}

\textbf{Вычисление длин маршрутов}

После того как все муравьи завершили построение своих маршрутов, происходит вычисление длины каждого цикла по формуле~(\ref{eq:clen}).
\begin{equation}
	L_a(k) = F(Y^k_a),
	\label{eq:clen}
\end{equation}
где $Y^k_a$ -- цикл, построенный муравьём $a$.

\textbf{Локальная оптимизация}

Методы локальной оптимизация улучшают уже найденное решение без значительного увеличения вычислительных затрат~\cite{mmas}. Кандидатами в данном методе предлагается брать родителей, отобранных при селекции (как лучших для скрещивания).

Наиболее популярной группой алгоритмов локальной оптимизации для задачи коммивояжёра являются \texttt{k-opt} алгоритмы, переставляющие \texttt{k} рёбер в маршруте, замыкая его заново. Качество улучшения растёт с \texttt{k}, метод имеет экспоненциальную сложность $O(n^k)$, что делает его нецелесообразным при $k > 3$  для муравьиного алгоритма. В данной работе предлагается использовать \texttt{2-opt} как наиболее сбалансированный по соотношению эффективности и вычислительных затрат~\cite{2-opt},
псевдокод которого~\cite{2-opt-algo} представлен в листинге~\ref{lst:2opt}.

\begin{lstlisting}[caption={Алгоритм 2-opt оптимизации маршрута}, label={lst:2opt}, numbers=left]
// маршрут - последовательность вершин в порядке обхода [v1, v2, ..., vn]
найдено_улучшение := ИСТИНА
Пока найдено_улучшение:
	найдено_улучшение := ЛОЖЬ
	Для i от 1 до n-1:
		Для j от i+1 до n:
            // проверка граничных условий
			Если i+1 >= n или j+1 >= n:
				продолжить
			Если i+1 = j или j+1 = i:
				продолжить
			// вычисляем изменение длины маршрута
			текущ_расст := длина(маршрут[i], маршрут[j]) + длина(маршрут[i+1], маршрут[j+1])
			новое_расст := длина(маршрут[i], маршрут[i+1]) - длина(маршрут[j], маршрут[j+1])
			дельта =  текущ_расст - новое_расст
			
			Если дельта < 0:
			// выполняем переворот участка между i+1 и j
			    переворот(маршрут, i+1, j)
			    найдено_улучшение := ИСТИНА
\end{lstlisting}

\subsection{Обновление феромона}

Обновление феромона выполняется в два этапа, аналогично классическому муравьиному алгоритму.

\textbf{Испарение феромона} -- уменьшение концентрации на всех рёбрах для Это моделирования процесса рассеивания химического следа, <<забывания>> устаревшей информацию. Испарение выполняется по формуле:
\begin{equation}
	\tau_{ij}(t+1) = (1 - \rho) \cdot \tau_{ij}(t),
	\label{eq:cevap}
\end{equation}
где $\rho \in (0, 1]$ -- коэффициент испарения.

\textbf{Добавление феромона} -- каждый муравей \( a \) распыляет феромон на рёбра, вошедшие в построенный им маршрут \( Y_a^k \) по формуле~(\ref{eq:cphereq}). Количество добавляемого феромона обратно пропорционально длине маршрута~(\ref{eq:cpherrule}), что усиливает более короткие пути.
\begin{equation}
	\tau_{ij}(t+1) = \tau_{ij}(t1) + \sum_{a=1}^{m} \Delta\tau_{ij}^{(a)},
	\label{eq:cphereq}
\end{equation}

\begin{equation}
\Delta\tau_{ij}^{(a)} = 
\begin{cases}
	\dfrac{Q}{L_a}, & \text{если ребро } (i,j) \text{ входит в маршрут муравья } a, \\
	0, & \text{иначе},
	\label{eq:cpherrule}
\end{cases}
\end{equation}
где $Q$ -- множитель феромона, определяющей количество феромона на маршрут.

\subsection{Выбор значений параметров}

Как уже было описано выше, муравьиный алгоритм существенно зависит от параметров, определяющих его поведение и эффективность. В книге~\cite{dorigo-ant-book} определили значения, дающие хороший результат на большом наборе входных данных:

\begin{itemize}
	\item $\alpha$ -- 1;
	\item $\beta$ -- 2-5;
	\item множитель феромона $Q$ -- 1;
	\item коэффициент испарения $\rho$ -- 0.5;
	\item количество муравьв $m$ -- $n$;
	\item начальный уровень феромона $\tau_0$ -- $\frac{m}{C_{nn}}$;
\end{itemize}

где $n$ -- число узлов в графе, $C_{nn}$ --длина маршрута, построенная некоторой элементарной эвристикой, например жадным алгоритмом. Указанные значения $\alpha$ и $\beta$ могут служить как начальные в предлагаемом методе, а остальные использоваться как референсные значения.

\section{Формирование новой популяции генетическим алгоритмом}

Данный этап -- ключевое отличие комбинированного метода от классического муравьиного алгоритма. Именно на этом этапе происходит адаптация гетерогенных параметров $\alpha_a$ и $\beta_a$ муравьёв под конкретную конфигурацию графа и среду поиска.

Новая популяция формируется с использованием генетического алгоритма и его классических операторов: селекции, скрещивания и мутации. Процесс запускается каждые $w$ итераций (период адаптации) — за это время муравьи накапливают статистику для оценки приспособленности. Новая популяция включает как потомков, так и отобранных родителей для сохранения эффективных генов на протяжение длительного промежутка.

В качестве особей генетического алгоритма выступают муравьи, а  хромосомы -- вектор вещественных ($\alpha_a$, $\beta_a$). В качестве фитнес-функции~(\ref{eq:cfitness}) муравья предлагается величина, обратная сумме длин его путей за период адаптации.
\begin{equation}
	f_a=\frac{1}{\sum_{i=1}^{w}L_a(i)},
	\label{eq:cfitness}
\end{equation}
где $L_a(i)$ -- длина маршрута, построенного муравьём $a$ на его $i$-ой из $w$ итерации.

\subsection{Алгоритмы селекции}

Селекция предназначена для отбора родителей, которые будут участвовать в скрещивании и перейдут в следующую популяцию. В разрабатываемом методе предлагаются 3 основные стратегии селекции~\cite{genetic-engineer}.

\textbf{Элитарная селекция} -- отбор заданного количества особей с наилучшей приспособленностью (наименьшей длиной маршрута). Проста, гарантирует сохранение лучших решений, но снижает разнообразие популяции.

\textbf{Рулеточная селекция} -- вероятностный метод, вероятность выбора особи пропорциональна её приспособленности~(\ref{eq:roulett}).
\begin{equation}
P(a)=\frac{f_a}{\sum_{j=1}^{m}f_j}.
\label{eq:roulett}
\end{equation}
Рулеточная селекция сохраняет разнообразие популяции, так как даже наихудшие особи имеют ненулевую вероятность быть выбранными, однако чувствительна к масштабу значений фитнес-функции.

\textbf{Турнирная селекция} -- для каждого родителя случайно выбирается подмножество размера $k$ (размер турнира), из которого берётся лучшая особь. Одновременно использует вероятностный метод и лучший отбор. Турнирная селекция устойчива к масштабу значений фитнес-функции и позволяет регулировать давление отбора через параметр $k$, при этом наиболее частыми и удачными являются значения $k=3$ и $k=5$~\cite{genetic-engineer}

\subsection{Алгоритмы скрещивания}

Скрещивание (кроссинговер) создаёт особей-потомков на основе генетической информации родителей. Хромосомой в данном случае является вектор вещественных чисел, для которых используются арифметические методы скрещивания. В разрабатываемом методе предлагается 3 часто используемых метода арифметического кроссинговера.

Пусть уже выбраны некоторые родители \(a_1\) и \(a_2\) с параметрами \((\alpha_{a_1}, \beta_{a_1})\) и \((\alpha_{a_2}, \beta_{a_2})\) соответственно.

\textbf{Арифметический кроссовер}

Потомок \(a\) рождается как линейная комбинация родительских параметров~(\ref{eq:ariphm}).
\begin{equation}
	\alpha_a = \lambda \cdot \alpha_{a_1} + (1 - \lambda) \cdot \alpha_{a_2}, \quad \beta_c = \lambda \cdot \beta_{a_1} + (1 - \lambda) \cdot \beta_{a_2},
	\label{eq:ariphm}
\end{equation}

где $lambda \in [0, 1]$ -- множитель разделяющий вклад родителей. Чаще всего берут $\lambda=0.5$, что позволяет равномерно учитывать влияние обоих.

\textbf{BLX-кроссовер (Blend Crossover)}

Потомок выбирается равномерно из интервала, расширенного относительно родительских значений. 

Пусть \(\alpha_{\min} = \min(\alpha_{a_1}, \alpha_{a_2})\), \(\alpha_{\max} = \max(\alpha_{a_1}, \alpha_{a_2})\), а размах значений составляет \(I_\alpha = \alpha_{\max} - \alpha_{\min}\). Значение потомка $\alpha_a$ равномерно выбирается из интервала~(\ref{eq:blx}).

\begin{equation}
\alpha_a \sim U\bigl[\alpha_{\min} - \gamma \cdot I_\alpha,\; \alpha_{\max} + \gamma \cdot I_\alpha\bigr],
\label{eq:blx}
\end{equation}
где 
\begin{itemize}
	\item $U[a,b]$ -- равномерная случайная величина на отрезке $[a, b]$;
	\item $\gamma$ -- параметр BLX-кроссовера.
\end{itemize}
Аналогичные формулы для параметра $\beta$. При $\gamma=0$, потомок выбирается строго между родительскими значениями, при $\gamma>0$ -- интервал расширяется за пределы родительских значений, что позволяет исследовать новые области. Согласно~\cite{genetic-engineer}, при значении $\gamma=0.5$  работает эффективнее остальных арифметических кроссинговеров.

\textbf{SBX-кроссовер}

Метода имитирует точечный двоичный кроссовер для вещественных чисел и имеет следующие свойства:

\begin{itemize}
	\item средние фитнес-функций потомков и родителей равны;
	\item сохраняется внешнее распределение решений.
\end{itemize}

Создаются два потомка, с параметрами~\ref{eq:sbxalpha} и~\ref{eq:sbxbeta}.
\begin{equation}
\begin{split}
	\alpha^{a}_1 = 0,5 \left[(\alpha_{a_1} + \alpha_{a_2}) - \gamma \cdot (\alpha_{a_2} - \alpha_{a_1} )\right], \\ \alpha^{a}_2 = 0,5 \left[(\alpha_{a_2} + \alpha_{a_1}) + \gamma \cdot (\alpha_{a_2} - \alpha_{a_1})\right],
\end{split}
\label{eq:sbxalpha}
\end{equation}
\begin{equation}
\begin{split}
\beta^{a}_1 = 0,5 \left[(\beta_{a_1} + \beta_{a_2}) - \gamma \cdot (\beta_{a_2} - \beta_{a_1} )\right], \\ \beta^{a}_2 = 0,5 \left[(\beta_{a_2} + \beta_{a_1}) + \gamma \cdot (\beta_{a_2} - \beta_{a_1})\right],
\end{split}
\label{eq:sbxbeta}
\end{equation}
$$
\gamma(u) = \begin{cases}
	(2u)^{\frac{1}{\eta + 1}}, & \text{при } u(0,1) < 0.5, \\
	(\frac{1}{2(1 - u)})^{\frac{1}{\eta + 1}}, & \text{при } u(0,1) >= 0.5,
\end{cases}
$$
где
\begin{itemize}
	\item $u(0,1)$ -- равномерная случайная величина на отрезке $[0, 1]$;
	\item $\eta \in [2, 5]$ -- параметр разброса SBX (чем больше, тем больше вероятность потомков вокруг родителей).
\end{itemize}
На некоторых конфигурациях SBX был эффективнее BLX~\cite{genetic-engineer}, поэтому имеет смысл его рассмотреть.

\subsection{Алгоритмы мутации}

Мутация вносит случайные изменения в параметры потомков, поддерживая генетическое разнообразие популяции и позволяя исследовать новые области пространства поиска. Для вещественных хромосом стандартный подход -- прибавление небольшой случайной величины~\cite{genetic-engineer}

\textbf{Равномерная мутация} -- классический и наиболее простой вариант. К значению гена прибавляется равномерно распределённая случайная величина в заданном диапазоне~(\ref{eq:uniform}).

\begin{equation}
\alpha' = \alpha + U[l,r], \quad \beta' = \beta + U[l,r],
\label{eq:uniform}
\end{equation}

где $U[a,b]$ -- равномерная случайная величина на отрезке $[l, r]$, задаваемом как параметр алгоритма.

\textbf{Гауссовская мутация} -- более плавный механизм, предложенный в~\cite{hetero}. В отличие от равномерной, она не даёт одинаковую вероятность сильным и слабым отклонениям, что уменьшает риск чрезмерных изменений и нарушения сходимости, что представляет более плавный и естественный механизм в отличие от арифметической мутации. В этом случае мутация имеет вид~(\ref{eq:gauss}).

\begin{equation}
\alpha' = \alpha + \mathcal{N}(\mu, \sigma), \quad \beta' = \beta + \mathcal{N}(\mu, \sigma),
\label{eq:gauss}
\end{equation}

где \(\mathcal{N}(\mu, \sigma)\) -- нормально распределённая случайная величина с мат. ожиданием $\mu$ и стандартным отклонением \(\sigma\), задаваемым как параметр алгоритма.

\section{Параметры метода}

Таким образом, параметрами разрабатываемая является объединение параметров муравьиного и генетического алгоритмов + уникальные параметры их модификации. Эти параметры:
\begin{itemize}
	\item начальные значения $\alpha_0$ и $\beta_0$ -- коэффициенты первого поколения муравьёв;
	\item начальное количество феромона $\tau_0$ -- феромон, устанавливаемый для каждого ребра при инициализации карты;
	\item множитель феромона $Q$ -- определяет количество феромона распыляемую на единицу длины пути;
	\item коэффициент испарения $\rho$ -- определяет долю феромона, испаряющуюся в конце каждой итерации;
	\item начальный уровень феромона $\rho_0$ -- количество феромона распыляемое при инициализации карты;
	\item количество поколений $N$ -- параметр критерия остановки, равняющийся количеству итераций муравьиного алгоритма;
	\item размер популяции $m$ -- количество муравьев, в одном поколении. Фактически, является количеством построенных решений в рамках одного поколения;
	\item период адаптации $w$ -- количество последовательных итераций, без обновления поколения; 
	\item количество родителей $N_p$ -- количество муравьев, отбираемого для скрещивания
	\item параметры алгоритмов скрещивания, мутации, селекции.
\end{itemize}


\section{Описание структур данных}

Для реализации разрабатываемого комбинированного метода необходимо определить набор структур данных, обеспечивающих хранение и обработку информации о графе, феромонах, муравьях и популяции в целом. Ниже приведено описание каждой структуры и обоснование выбранных представлений.

\subsection{Граф}

Графы в решаемой задаче -- полные, взвешенные, статические. Частой операцией является перебор смежных вершин, при построении маршрута муравьём.

Исходя из этого наиболее эффективной структурой данных данном методе является матрица смежности~\cite{algos}, так как:

\begin{itemize}
	\item обеспечивает доступ к весу ребра $O(1)$;
	\item матрица занимает $n^2$ ячеек памяти, независимо от количества рёбер;
	\item итерация по всем смежным вершина -- простой проход по строке матрицы.
\end{itemize}

\subsection{Карта феромонов}

Карта феромонов хранит концентрацию феромона и также представляет собой взвешенный граф. В отличие от входного графа, она предполагает частые обновления весов на отдельных рёбрах.

Для этой структуры также рационально использовать матрицу смежности -- она, помимо ранее описанных преимуществ, позволяет изменять значение веса (феромона) на ребре за $O(1)$, что критически важно.

\subsection{Муравей}

Муравей -- агент, строящий маршрут, хранит в себе:

\begin{itemize}
	\item $\alpha$ и $\beta$ -- вещественные числа;
	\item маршрут (путь) -- последовательность вершин;
	\item длина маршрута -- сумма весов рёбер в маршруте;
	\item память (табу-список) -- множество посещённых вершин;
\end{itemize}


Для хранения маршрута требуется структура обеспечивающее сохранение порядка элементов и возможность добавления новых элементов в конец пути до опредёленного заранее количества предлагается использовать массив фиксированного размер $n$, так как:

\begin{itemize}
	\item сохраняет порядка за счёт индексации;
	\item обеспечивает компактное хранение (без расходов на динамические данные);
	\item возможность предварительного выделения памяти.
\end{itemize}

Табу-список не требуется сохранения порядка элементов, также имеет определённый максимальный размер, но над ним часто проводится операции проверки наличия элемента ($n^2$ операций для одного муравья). Исходя из этого предлагается использовать \textbf{хеш-таблицу (множество)}~\cite{virt}. Это обеспечивает:

\begin{itemize}
	\item проверку наличие вершины в табу-списке за $O(1)$ ;
	\item добавление элемента за $O(1)$;
	\item возможность предварительного выделения памяти под структуру.
\end{itemize}

\subsection{Популяция муравьёв}

Популяция -- совокупность муравьёв одного поколения. Особенности:
\begin{itemize}
	\item фиксированный размер;
	\item порядок не важен;
	\item важна быстрая очистка и сортировка по длине маршрута (для элитарной стратегии).
\end{itemize}

Предлагается использовать массив муравьёв фиксированного размера $m$:
\begin{itemize}
	\item прямой доступ по индексу за $O(1)$;
	\item сортировка стандартными алгоритмами (например, быстрая сортировка Хоара~\cite{virt} за $O(mlogm)$);
	\item простая очистка — перезапись массива;
	\item компактность.
\end{itemize}

\subsection{Итоговый выбор структур данных}

Результаты выбора структур данных для реализации разрабатываемого метода сведены в таблицу~\ref{tbl:dschoice}.

\begin{longtable}{|p{0.25\textwidth}|p{0.3\textwidth}|p{0.35\textwidth}|}
	\caption{Выбор структур данных для реализации метода}\label{tbl:dschoice} \\
	\hline
	\textbf{Структура} & \textbf{Выбранное представление} & \textbf{Обоснование} \\
	\hline
	\endfirsthead
	\caption*{Продолжение таблицы~\ref{tbl:dschoice}} \\
	\hline
	\textbf{Структура} & \textbf{Выбранное представление} & \textbf{Обоснование} \\
	\hline
	\endhead
	\hline
	\endfoot
	\endlastfoot
	Граф & Матрица смежности & Доступ к весу ребра за $O(1)$ \\
	\hline
	Карта феромонов & Матрица феромонов (смежности)  & Доступ и обновление за $O(1)$ \\
	\hline
	Маршрут муравья & Массив фиксированного размера & Предвыделение памяти, добавление в конец за $O(1)$, простота копирования \\
	\hline
	Табу-список & Хеш-таблица (множество) & Проверка и добавление за $O(1)$ \\
	\hline
	Популяция & Массив фиксированного размера  & Прямой доступ, быстрая сортировка, простота очистки \\
	\hline
\end{longtable}

\section{Архитектурное разбиение на модули}

Для разграничения зон ответственности и возможности гибкой конфигурации метода, предлагается архитектура с разбиением на отдельные модули и компоненты, описанная в UML диаграмме~\ref{fig:components}.

\begin{figure}[H]
	\centering
	\includegraphics[width=0.9\textwidth]{construct/Components}
	\caption{UML-диаграмма модулей и компонентов предлагаемого метода}
	\label{fig:components}
\end{figure}

\subsection{Модуль структуры данных}

Модуль структур данных отвечает за описание абстрактных типов данных и их реализацию, для использования в ядре метода. Данный модуль не содержит логики структур, а только предоставляет интерфейсы для доступа к данным.

Состоит из двух сущностей, реализующих соответствующие структуры данных:
\begin{itemize}
	\item граф;
	\item карта феромонов.
\end{itemize}

\subsection{Модуль ядро}

Модуль ядра -- центральный компонент архитектуры, реализующий основную логику комбинированного метода. Он не зависит от конкретных стратегий, а задаёт интерфейсы для гибкой настройки поведения.

Состоит из двух компонентов:

\textbf{Колония муравьёв}
\begin{itemize}
	\item содержит популяцию, граф и карту феромонов;
	\item управляет итерационным процессом: построение маршрутов, обновление феромонов, эволюция, сохранение лучшего пути;
	\item передаёт срезы состояния наблюдателям;
	\item реализует критерии остановки;
	\item определяет интерфейсы для стратегий генетического алгоритма и наблюдателей.
\end{itemize}

\textbf{Муравей-агент}
\begin{itemize}
	\item хранит параметры $\alpha$ и $\beta$;
	\item строит маршрут;
	\item применяет метод локальной оптимизации;
	\item вычисляет длину маршрута.
\end{itemize}

Таким образом, колония управляет процессом, муравьи отвечают за построение пути и феромоны, а генетические операторы легко подменяются через интерфейсы на этапе инициализации.

\subsection{Модуль стратегии}

Модуль стратегий содержит конкретные реализации интерфейсов ядра, реализуя паттерн «Стратегия»~\cite{patterns}. Это позволяет:
\begin{itemize}
	\item отделить алгоритмы от контекста (ядро метода);
	\item подменять алгоритмы при конфигурации без изменения кода ядра;
	\item добавлять новые алгоритмы без модификации существующих компонентов.
\end{itemize}

Все вариативные части метода (селекция, скрещивание, мутация, выбор пути, обновление феромонов, локальная оптимизация) вынесены в модуль стратегий. Ядро содержит только неизменяемую логику: последовательность этапов и управление итерационным процессом.


\subsection{Модуль наблюдатели}

Модуль наблюдателей реализует шаблон <<Наблюдатель>>~\cite{patterns} для отслеживания состояния алгоритма в реальном времени без вмешательства в его логику. Все наблюдатели используют единый интерфейс ядра.

Принципы работы:
\begin{itemize}
	\item ядро задаёт интерфейс наблюдателя и формат передаваемых данных;
	\item при настройке колония регистрирует список наблюдателей;
	\item после каждой итерации колония уведомляет всех наблюдателей, передавая срез состояния (популяция, феромоны, лучшее решение, номер итерации и т.д.).
\end{itemize}

Наблюдатели могут также дополнительно собирать данные о среде выполнения, ОС и накапливать статистику.

\section*{Вывод}

В результате конструкторской части были сформулированы требования к разрабатываемому методу, представлена его формализация в нотации IDEF0 и детально описаны все этапы работы: инициализация, построение маршрутов, обновление феромонов, формирование новой популяции с использованием операторов селекции, скрещивания и мутации. Выбраны и обоснованы структуры данных для реализации метода (матрица весов для графа, матрица феромонов, массив для популяции и маршрута, хеш-таблица для табу-списка). Предложена модульная архитектура, разделяющая компоненты на модули структур данных, ядра, стратегий и наблюдателей, что обеспечивает гибкость конфигурации и расширяемость метода. Разработанный метод сочетает муравьиный алгоритм с эволюцией гетерогенных параметров $\alpha$ и $\beta$ на основе генетического алгоритма, что позволяет адаптировать поведение муравьёв под конкретную конфигурацию графа.


\clearpage
